
\chapter{三角}


%\section{正弦、余弦、正切、余切}


\section{锐角的正弦、余弦、正切、余切}
\zsd 锐角的正弦、余弦、正切、余切

A组：
\begin{enumerate}
\item 填空\\
\begin{minipage}{\textwidth}
\centering
\includegraphics[scale=0.9]{./bodyxb/src/main42xbb2ch06ruijiandesanjiaobi01t1.png}
\end{minipage}
\begin{enumerate}[(1)]
\item $\sin{A}=$\blkda{$\frac{a}{c}$}, $\cos{A}=$\blkda{$\frac{b}{c}$}, $\tan{A}=$\blkda{$\frac{a}{b}$}, $\cot{A}=$\blkda{$\frac{b}{a}$}. 
\item $\sin^2A+\cos^2A=$\blkda{1}. 
\item $\sin(\ang{90}-A)=$\blkda{$cosA$}, $\cos(\ang{90}-A)=$\blkda{$sinA$}, \\
$\tan(\ang{90}-A)=$\blkda{$cotA$}, $\cot(\ang{90}-A)=$\blkda{$tanA$}. 
\item 
\begin{minipage}[t]{\textwidth}
\begin{tblr}{hlines,vlines,column{1}={2cm,c},column{2,3,4,5}={2.5cm,c},row{1}={blue!10}}
角度$\alpha$ & $\sin\alpha$ & $\cos\alpha$ & $\tan\alpha$ & $\cot\alpha$\\
$\ang{30}$ & \blkda{$\frac{1}{2}$} & \blkda{$\frac{\sqrt{3}}{2}$} &\blkda{$\frac{\sqrt{3}}{3}$} & \blkda{$\sqrt{3}$} \\
$\ang{45}$ & \blkda{$\frac{\sqrt{2}}{2}$} & \blkda{$\frac{\sqrt{2}}{2}$} &\blkda{$1$} & \blkda{$1$} \\
$\ang{60}$ & \blkda{$\frac{\sqrt{3}}{2}$} & \blkda{$\frac{1}{2}$} & \blkda{$\sqrt{3}$} & \blkda{$\frac{\sqrt{3}}{3}$} \\
\end{tblr}
\end{minipage}
\end{enumerate}

\item  如图, 在 $x$ 轴的上方, 直角 $\angle B O A$ 绕原点 $O$ 按顺时针方向旋转. 若 $\angle \mathrm{BOA}$ 的两边分别与函数 $y=-\frac{1}{x}, y=\frac{2}{x}$ 的图象交于$B 、 A$ 两点, 则 $\angle O A B$ 大小的变化趋势为\dotfill(\kh{D})
\xx{逐渐变小}{逐渐变大}{时大时小}{保持不变}
\begin{minipage}{\textwidth}
\centering
\includegraphics[scale=0.9]{./bodyxb/src/31247891329857981275.jpg}
\end{minipage}
\begin{jieda}
    过$A, B$做$x$轴垂线，设垂足为$A', B'$，根据三角形相似可知$|OB'| \cdot |AA'| = \sqrt{2}$，进一步可以推导出$\dfrac{|OA|}{|OB|} = \sqrt{2}$，即$\angle \mathrm{BOA}$的正切值固定，因此选D.
\end{jieda}


\item 试判断下列命题是否正确:
\begin{enumerate}
    \item 小于 $90^{\circ}$ 的角是锐角;\dotfill()
    \item 第一象限角小于第二象限角；\dotfill()
    \item 终边相同的角一定相等；\dotfill()
    \item 相等的角终边一定相同；\dotfill()
    \item 若 $90^{\circ}<a<180^{\circ}$, 则 $a$ 是第二象限角.\dotfill()
\end{enumerate}
\begin{jieda}
\begin{enumerate}[(1)]
\item $-\ang{45}$角不是锐角, 故假. 
\item $\ang{30}>\ang{120}-\ang{360}=-\ang{240}$, 故假. 
\item $\alpha=\ang{30}$, $\beta=\ang{30}+\ang{360}$, 故假. 
\item 真. 如果有破绽的话, 破绽在于始边是否重合. 
\item 真. 
\end{enumerate}
\end{jieda}



\item $-2021^{\circ}$ 是第\blkda{II}象限角.
\begin{jieda}
$-2021=-5\cdot360-221=-6\cdot 360+139$, 故II. 
\end{jieda}


\item 角 $a$ 是第四象限角, 则角 $180^{\circ}-a$ 是第\blkda{三}象限角.


\item 分别写出第一、二、三、四象限角的集合. 
\begin{jieda}
一: $\{\theta\mid\theta=k\cdot\ang{360}+\alpha,\ang{0}\alpha<\ang{90},k\zz\}$; \\
二: $\{\theta\mid\theta=k\cdot\ang{360}+\alpha,\ang{90}\alpha<\ang{180},k\zz\}$; \\
三: $\{\theta\mid\theta=k\cdot\ang{360}+\alpha,\ang{180}\alpha<\ang{270},k\zz\}$; \\
四: $\{\theta\mid\theta=k\cdot\ang{360}+\alpha,\ang{270}\alpha<\ang{360},k\zz\}$. 
\end{jieda}


\end{enumerate}

B组：
\begin{enumerate}
\item $\dfrac{\pi}{5}$=\blkda{36}度; $-72^{\circ}$=\blkda{$-\dfrac{2\pi}{5}$}弧度.


\item 若时钟的分针转了2小时15分, 则所转的角度是\blkda{$-\ang{810}$}. 
\begin{jieda}
$s=vt=360\cdot 2.25=810$, 就是这个东西. 
\end{jieda}


\item 角的顶点在坐标系的原点, 始边与$x$轴的正半轴重合, 那么终边在下列位置的角的集合分别是: 
\begin{enumerate}[(1)]
\item $x$轴正半轴: \blkda[5]{$\{\theta\mid \theta=k\cdot\ang{360},k\zz\}$}; 
\item $x$轴: \blkda[5]{$\{\theta\mid \theta=k\cdot\ang{180},k\zz\}$}; 
\item $y$轴负半轴: \blkda[5]{$\{\theta\mid \theta=k\cdot\ang{360}-\ang{90},k\zz\}$}; 
\item $y$轴: \blkda[5]{$\{\theta\mid \theta=k\cdot\ang{180}+\ang{90},k\zz\}$}; 
% \item 坐标轴: \blkda{}; 
\item 直线$y=-x$: \blkda[5]{$\{\theta\mid \theta=k\cdot\ang{180}+\ang{135},k\zz\}$}; 
\item 直线$y=\pm{x}$: \blkda[5]{$\{\theta\mid \theta=k\cdot\ang{90}+\ang{45},k\zz\}$}. 
\end{enumerate}
\begin{jieda}
以(6)为例: 每逆时针旋转$\ang{90}$, 恰好到达一条射线. 于是从角的角度看, 遍历了所有的角. 
相当于有了一个起始值$\ang{45}$, 每次加上或减去$\ang{90}$, 就有了所有角. 
\end{jieda}



\item 已知$\alpha$与$\beta$都是锐角, $\alpha+\beta$的终边与$-280^{\circ}$的终边相同; 
$\alpha-\beta$的终边与$-670^{\circ}$的终边相同, 则$\alpha$=\blkda{$\ang{65}$}, $\beta$=\blkda{$\ang{15}$}. 
\begin{jieda}
$\begin{cases}
\alpha+\beta=m\cdot\ang{360}-\ang{280}\\
\alpha-\beta=n\cdot\ang{360}-\ang{670}
\end{cases}$, 
故$\alpha=(m+n)\cdot\ang{180}-\ang{475}$, $\beta=(m-n)\cdot\ang{180}+\ang{195}$. \\
又$\alpha,\beta$是锐角, 故$\alpha=\ang{65}$, $\beta=\ang{15}$. 
\end{jieda}


\item 已知圆周上有一点$P$(初始位置在$x$轴正半轴上)依逆时针方向作匀速圆周运动. 已知点$P$每分钟转过$\theta$角($\ang{0}<\theta\leq\ang{180})$, 
2分钟到达第三象限, 15分钟回到原来的位置, 则角$\theta$等于\blkda[2]{$\ang{96}$或$\ang{120}$}. 
\begin{jieda}
$\begin{cases}
\ang{0}<\theta<\ang{180}\\
\ang{180}<2\theta<\ang{270}\\
15\theta=k\cdot\ang{360}
\end{cases}$, 
将第三式代入前两式得到$k$的范围, 结合$k\zz$得$k=4,5$, 故$\theta=\ang{96}$或$\ang{120}$. 
\end{jieda}


\item 若$\alpha$与$\beta$是终边相同的角, 则$\alpha-\beta$的终边一定在\dotfill(\kh{A})
\xx{$x$轴的非负半轴上}
{$x$轴的非正半轴上}
{$y$轴的非正半轴上}
{$y$轴的非负半轴上}



\item 设$k\zz$, 终边相同的一组角是\dotfill(\kh{C})
\xx{$k\cdot 90^{\circ}$与$k\cdot 180^{\circ}+90^{\circ}$}
{$k\cdot 180^{\circ}+60^{\circ}$与$k\cdot 60^{\circ}$}
{$(2k-1)\cdot 180^{\circ}$与$(4k\pm 1)\cdot 180^{\circ}$}
{$k\cdot 180^{\circ}+30^{\circ}$与$k\cdot 360^{\circ}-150^{\circ}$}



\item 在与2021弧度角的终边相同的角中, 绝对值最小的角是\blkda[3]{$2021-644\pi$}. 
\begin{jieda}
$\abs{2021+2k\pi}_{\min}=\abs{2021-322\cdot2\pi}$. 
\end{jieda}



\item 若角$\alpha$和角$\beta$的和是$1$弧度, 差为$1^{\circ}$, 则$\alpha$和$\beta$的弧度数分别为\blkda{$\frac{1}{2}\pm\frac{\pi}{360}$}. 

\item 若中心角为$1^{\circ}$的扇形, 它的弧长为$2\pi$, 则该扇形所在圆的半径为\blkda{360}. 
\begin{jieda}
$r=\dfrac{l}{\alpha}=\dfrac{2\pi}{\frac{\pi}{180}}=360$. 
\end{jieda}


\item 若扇形的圆心角是$120^{\circ}$, 则该扇形面积与其内切圆面积的比值是\blkda{$\dfrac{7+4\sqrt{3}}{9}$}. 
\begin{jieda}
相切找切点, 于是结合简单的几何知识易得$R=r+\dfrac{2r}{\sqrt{3}}$. 
\end{jieda}


\item 若一条弦长等于半径的$\sqrt{2}$倍, 则此弦所对圆心角等于\blkda{$\dfrac{\pi}{2}$}弧度.
\begin{jieda}
$\sin\frac{\theta}{2}=\frac{\sqrt{2}}{2}$. 
\end{jieda}

\item 如果$1$弧度的圆心角所对的弦长为$2$, 那么这个圆心角所对的弧长是\blkda{$\dfrac{1}{\sin\frac{1}{2}}$}.
\begin{jieda}
``弧''与``弦''; 弦取中点. 
\end{jieda}



\item 若扇形的周长是$16$, 圆心角是$2$弧度, 则扇形面积是\blkda{$16$}.


\item
\begin{enumerate}[(1)]
\item 若角$\alpha$的终边过点$(2^x,\log_{\frac{1}{2}}2)$, $x\zr$, 则$\dfrac{\alpha}{2}$终边在第\blkda{II、IV}象限. 

\item 设$\alpha$是第二象限角, 则$\dfrac{\alpha}{3}$的终边不在\blkda{第三象限}.

\end{enumerate}

\begin{jieda}

先看图: \\
\begin{minipage}{\textwidth}
\centering
\includegraphics[scale=0.8]{./bodyxb/src/main42chap05renyijiaodesanjiaobi01.pdf}
\figcaption{半角、倍角与象限}
\end{minipage}
该图的信息有点多: 

\begin{itemize}
\item 由$\alpha$所在的象限, 可得$\dfrac{\alpha}{2}$所在象限, 且只能是象限中的一部分. 
\item 由$\alpha$所在的象限, 可得$2\alpha$所在象限. 
\item 可推广到$\dfrac{\alpha}{n}(n\zn)$的情形. 
\end{itemize}

\begin{enumerate}[(1)]
\item $2^x>0, \log_{\frac{1}{2}}2=-1$, 故$\alpha$是第四象限角, 于是$\dfrac{\alpha}{2}$是第II,IV象限角. 

\item 每个象限等分为3份, 然后逆时针依次标上1234(周期). 

\end{enumerate}

\end{jieda}



\item 一扇形周长是\SI{32}{cm}, 扇形的圆心角为多少弧度时, 这个扇形的面积最大? 最大面积是多少? 
\begin{jieda}
$l+2r=32$, $S=\frac{1}{2}lr=\frac{1}{4}\cdot l\cdot2r\leq\frac{1}{4}(\frac{l+2r}{2})^2=64$. \\
$\alpha=\frac{l}{r}=\SI{2}{rad}$, $S_{\min}=\SI{64}{cm^2}$. 
\end{jieda}


\item 已知扇形的圆心角为 $\frac{\pi}{3}$, 半径为 $R$ , 圆 $C$ 与扇形的两条半径及扇形的弧都相切, 求圆 $C $中圆心角为 $\frac{\pi}{3}$ 的扇形与原扇形的面积之比.
\begin{jieda}
    根据几何关系，半径之比为$\dfrac{1}{3}$，因此面积之比$\dfrac{1}{9}
\end{jieda}

\end{enumerate}

C组：
\begin{enumerate}
    \item 求$\dfrac{1-2log_65 \cdot log_{10}3 \cdot log_{15}2}{log_65 \cdot log_{10}3 + log_{10}3 \cdot log_{15}2 + log_{15}2 \cdot log_65}$的值.
    \begin{jieda}
        考虑换底，$log_65 = \dfrac{log_25}{log_23+1}, log_{10}3 = \dfrac{log_23}{log_25+1}, log_{15}2 = \dfrac{1}{log_23 + log_25}$\\
        故设$a = log_23, b = log_25$\\
        原式等于$\dfrac{1 - 2 \cdot \dfrac{b}{a+1} \cdot \dfrac{a}{b+1} \cdot \dfrac{1}{a+b}}{\dfrac{b}{a+1}\cdot \dfrac{a}{b+1} + \dfrac{a}{b+1} \cdot \dfrac{1}{a+b} + \dfrac{1}{a+b} \cdot \dfrac{b}{a+1}} = \dfrac{(a+1)(b+1)(a+b) - 2ab}{ab(a+b) + a(a+1) + b(b+1)} = 1$
    \end{jieda}
    \item 已知定义在$R^{+}$上的函数$f(x) = \left \{\begin{array}{ll}
        |log_3x - 1|, & 0 < x \leq 9 \\
        4-\sqrt{x}, & x > 9
    \end{array} \right.$，设$a, b, c$是三个互不相同的实数，满足$f(a) = f(b) = f(c)$，求$abc$的取值范围.
    \begin{jieda}
        不妨设$a < b < c$，易知$0 < a < 3 < b < 9< c$，因此有$1 - log_3a = log_3b - 1$，即$ab = 9$. \\
        故$f(a) \in (0, 1)$，因此$c \in (9, 16)$，因此有$abc \in (81, 144)$
    \end{jieda}
    \item 若关于$x$的方程$\dfrac{|x|}{x+6} = kx^2$有四个不同的实数解，求实数$k$的取值范围.
    \begin{jieda}
        易知必有一根为0 \\
        参变分离有$\dfrac{1}{k} = \dfrac{(x+6)x^2}{|x|} = \left \{ \begin{array}{ll}
            \dfrac{(x+6)x^2}{x}, & x > 0 \\
            -\dfrac{(x+6)x^2}{x},  & x < 0
        \end{array}\right.$
        由于原方程必有一根为0，故可以转化为方程$\dfrac{1}{k} = \left \{ \begin{array}{ll}
            (x+6)x, & x > 0 \\
            -(x+6)x,  & x < 0
        \end{array}\right.$有三个不等实根.因此$\dfrac{1}{k} \in (0, 9)$，即$k \in (\dfrac{1}{9}, +\infty)$
    \end{jieda}
\end{enumerate}

\section{任意角的正弦、余弦、正切、余切（1）}
\zsd 任意角的正弦、余弦、正切、余切及其符号. 同角三角比关系（求值）

A组：
\begin{enumerate}
\item 已知角$\alpha$的终边经过下列各点, 求角$\alpha$的正弦、余弦、正切、余切. 
\begin{enumerate}[(1)]
\item $A(-8,-6)$,     则$\sin\alpha$=\blkda{$-\frac{3}{5}$}, $\cos\alpha$=\blkda{$-\frac{4}{5}$}, $\tan\alpha$=\blkda{$\frac{3}{4}$}, $\cot\alpha$=\blkda{$\frac{4}{3}$}; 
\item $B(\sqrt3,-1)$, 则$\sin\alpha$=\blkda{$-\frac{1}{2}$}, $\cos\alpha$=\blkda{$\frac{\sqrt{3}}{2}$}, $\tan\alpha$=\blkda{$-\frac{\sqrt{3}}{3}$}, $\cot\alpha$=\blkda{$-\sqrt{3}$}. 
\end{enumerate}


\item 下列各式为正号的是\dotfill(\kh{C})
\xx{$\cos2-\sin2$}
{$\cos2\cdot\sin2$}
{$\tan2\cdot\sec2$}
{$\sin2\cdot\tan2$}
\begin{jieda}
$2\in(\frac{\pi}{2},\pi)$, 故$\sin2>0$, $\cos2<0$, $\tan2<0$, $\cot2<0$, $\sec2<0$, $\csc2>0$. 
\end{jieda}



\item 填表:

\begin{tabular}{|c|c|c|c|}\hline
$\alpha$ & $\sin\alpha$ & $\cos\alpha$ & $\tan\alpha$\\\hline
$\frac{2\pi}{3}$ & \blkda{$\frac{\sqrt{3}}{2}$} & \blkda{$-\frac{1}{2}$} & \blkda{$-\sqrt{3}$} \\\hline
$\frac{5\pi}{4}$ & \blkda{$-\frac{\sqrt{2}}{2}$} & \blkda{$-\frac{\sqrt{2}}{2}$} & \blkda{1} \\\hline
$\frac{5\pi}{6}$ & \blkda{$\frac{1}{2}$} & \blkda{$-\frac{\sqrt{3}}{2}$} & \blkda{$-\frac{\sqrt{3}}{3}$}\\\hline
\end{tabular}


\item 填表:

\begin{tabular}{|c|c|c|c|c|c|}\hline
$\alpha$ & 点$P$的坐标 & ~~$\sin\alpha$~~ & ~~$\cos\alpha$~~ & ~~$\tan\alpha$~~ & ~~$\cot\alpha$~~ \\\hline
$\alpha=2k\pi(k\zz)$ & $(1,0)$ & 0 & 1 & 0 & 不存在\\\hline
\blkda{$\alpha=2k\pi+\frac{\pi}{2}(k\zz)$} & $(0,1)$ & \blkda{1} & \blkda{0} & \blkda{不存在} &  \blkda{0} \\\hline
\blkda{$\alpha=2k\pi+\pi(k\zz)$} & $(-1,0)$  & \blkda{0} & \blkda{-1} & \blkda{0} & \blkda{不存在} \\\hline
\blkda{$\alpha=2k\pi-\frac{\pi}{2}(k\zz)$}& $(0,-1)$ & \blkda{-1} & \blkda{0} & \blkda{不存在}  & \blkda{0}\\\hline
\end{tabular}



\item 已知$\sin\alpha=\dfrac{4}{5}$, $\alpha\in(0,\pi)$, 则$\tan\alpha$等于\blkda{$\pm\dfrac{4}{3}$}. 
\begin{jieda}
三角比的值分为两部分: 符号与绝对值. 结合定义, 直接赋值: $y=4,r=5$, 于是$x=\pm3$. 
\end{jieda}


\item 已知角 $\theta$ 满足 $\sin \theta<0$ 且 $\cos \theta>0$, 则角 $\theta$ 是第\blkda{四}象限的角.

\item 已知角 $a$ 的终边过点 $P(2 a, 3 a)(a<0)$, 则角 $a$ 的余弦值为\blkda{$-\frac{2\sqrt{13}}{13}$}.


\item 已知角 $\theta$ 满足 $\frac{\sin \theta}{\cos \theta}>0$, 则角 $\theta$ 是第\blkda{一或三}象限的角.


\item 设 $x \in \mathrm{R}$, 且 $x \neq k \pi+\frac{\pi}{2}(k \in Z)$, 则 $\frac{\sin x}{|\sin x|}+\frac{|\cos x|}{\cos x}+\frac{\tan x}{|\tan x|}$ 的取值集合为\blkda{$\{-1, 3\}$}


\end{enumerate}

B组：
\begin{enumerate}
\item 计算: $\cos\dfrac{\pi}{3}-\tan\dfrac{\pi}{4}+\dfrac{3}{4}\tan^2\dfrac{7\pi}{6}-\sin\dfrac{5\pi}{6}+\cos^2\dfrac{\pi}{6}+\sin\dfrac{3\pi}{2}$=\blkda{$-1$}. 

\item 若$\sin\theta=\dfrac{m-3}{m+5}$, $\cos\theta=\dfrac{4-2m}{m+5}$都有意义, 则$m$=\blkda{0或8}. 
\begin{jieda}
$\sin^2\theta+\cos^2\theta=1$. 
\end{jieda}


\item 已知 $\alpha$ 为第二象限的角, 其终边上有一点 $P(x, \sqrt{5})$, 且 $\cos \alpha=\frac{\sqrt{2}}{4} x$, 求 $\tan \alpha =$ \blkda{$-\frac{\sqrt{15}}{3}$}.
\begin{jieda}
    $\dfrac{1}{r} = \dfrac{\sqrt{2}}{4}$，因此$sin\alpha = \dfrac{\sqrt{5}}{r} = \sqrt{5} \cdot \dfrac{\sqrt{2}}{4} = \dfrac{\sqrt{10}}{4}$，故$cos\alpha = -\sqrt{1 - \dfrac{10}{16}} = -\dfrac{\sqrt{6}}{4}$，因此$tan\alpha = \dfrac{sin\alpha}{cos\alpha} = -\dfrac{\sqrt{10}}{\sqrt{6}} = -\dfrac{\sqrt{15}}{3}$
\end{jieda}


\item 已知角$\theta$的终边上有一点$P(4,y)$, 且$\sin\theta=-\dfrac{2\sqrt{5}}{5}$, 则$y$=\blkda{$-8$}. 

\item 已知角$\alpha$的终边与射线$5x+12y=0(x\leq 0)$的图像重合, $\cos\alpha+\cot\alpha-\tan\alpha$的值为\blkda{$-\dfrac{2267}{780}$}. 
\begin{jieda}
$5x+12y=0(x\leq 0)$上取点$(-12,5)$, 
则$E=\dfrac{-12}{13}+\dfrac{-12}{5}-\dfrac{5}{-12}=-\dfrac{2267}{780}$. 
\end{jieda}


\item 已知集合$A=\left\{y\mid y=\dfrac{\abs{\sin{x}}}{\sin{x}}+\dfrac{\cos{x}}{\abs{\cos{x}}}+\dfrac{\abs{\tan{x}}}{\tan{x}}\right\}$, 
用列举法表示集合$A$是\blkda{$\{-1,3\}$}. 
\begin{jieda}
四个象限讨论一下即可. 
\end{jieda}


\item 若角$\alpha$是第三象限角, 则
\ding{172} $\sin\alpha+\cos\alpha<0$; 
\ding{173} $\tan\alpha-\sin\alpha>0$; 
\ding{174} $\cot\alpha\cdot\csc\alpha<0$; 
\ding{175} $\sin\alpha\cdot\sec\alpha>0$. 
其中正确的是\blkda{\ding{172}\ding{173}\ding{174}\ding{175}}.  

\item 已知角$\alpha$的终边过点$(3a-9,a+2)$, 且$\cos\alpha\leq 0$, $\sin\alpha>0$, 则$a$的取值范围是\blkda{$(-2,3]$}. 

\item 在$\triangle{ABC}$中, 有$\cos{A}\cdot\tan{B}\cdot\cot{C}<0$, 则这三角形是\blkda{钝角}三角形.
\begin{jieda}
首先没有直角, 否则要么积为零, 要么积不存在. \\
其次$\cos,\tan,\cot$在$(0,\frac{\pi}{2})$上都是正的, 在$(\frac{\pi}{2},\pi)$上都是负的, \\
故结合条件可得$\cos{A},\tan{B},\cot{C}$中一负两正, 于是钝角三角形. 
\end{jieda}



\item 设$y=\sin\alpha\cdot\cos\alpha\cdot\cot\alpha$, 且$\alpha$是象限角, 则$y$的符号为\dotfill(\kh{A})
\xx{恒正}{恒负}{可能为0}{不确定}



\item 若角$\alpha$的终边落在直线$x+y=0$上, 则$\dfrac{\sin\alpha}{\sqrt{1-\sin^2\alpha}}+\dfrac{\sqrt{1-\cos^2\alpha}}{\cos\alpha}$的值等于\blkda{$0$}.
\begin{jieda}
考虑一下符号基本上就可以把答案选出来了. 
\end{jieda}


\item 已知实数$\alpha,\beta$满足$\abs{\cos\alpha-\cos\beta}=\abs{\cos\alpha}+\abs{\cos\beta}$, 且$\alpha\in(\frac{\pi}{2},\pi)$, 
则化简$\sqrt{(\cos\alpha-\cos\beta)^2}$的结果是\dotfill(\kh{C})
\xx{$\cos\alpha-\cos\beta$}
{$\abs{\cos\alpha}-\abs{\cos\beta}$}
{$\cos\beta-\cos\alpha$}
{$\abs{\cos\alpha}+\abs{\cos\beta}$}
\begin{jieda}
$\abs{\cos\alpha-\cos\beta}=\abs{\cos\alpha}+\abs{\cos\beta}\iff\cos\alpha\cos\beta\leq 0$. \\
又$\cos\alpha<0$, 故$\cos\beta\geq 0$. \\
故$E=\abs{\cos\alpha-\cos\beta}=\cos\beta-\cos\alpha$. 
\end{jieda}


\item 已知角$\alpha$的终边上一点$P$与点$A(-3,2)$关于$y$轴对称, 角$\beta$的终边上一点$Q$与点$A$关于原点对称, 求$2\sin\alpha+3\sin\beta$的值. 
\begin{jieda}
$A(-3,2)\xleftrightarrow{y\text{轴}}P(3,2)$, $A(-3,2)\xleftrightarrow{\text{原点}}Q(3,-2)$. \\
故$2\sin\alpha+3\sin\beta=2\cdot\dfrac{2}{\sqrt{13}}+3\cdot\dfrac{-2}{\sqrt{13}}=-\dfrac{2\sqrt{13}}{13}$. 
\end{jieda}


\item 已知角$\theta$终边上一点$P(-\sqrt{3},m)$, 且$\sin\theta=\dfrac{\sqrt{2}m}{4}$, 求$\cos\theta$和$\tan\theta$的值. 
\begin{jieda}
$\dfrac{m}{\sqrt{3+m^2}}=\dfrac{\sqrt{2}m}{4}$, 则$m=0,\pm\sqrt{5}$. 
\begin{itemize}
\item $m=0$时, $\cos\theta=-1$, $\tan\theta=0$; 
\item $m=\sqrt{5}$时, $\cos\theta=-\dfrac{\sqrt{6}}{4}$, $\tan\theta=-\dfrac{\sqrt{15}}{3}$; 
\item $m=-\sqrt{5}$时, $\cos\theta=-\dfrac{\sqrt{6}}{4}$, $\tan\theta=\dfrac{\sqrt{15}}{3}$. 
\end{itemize}
\end{jieda}


\item 已知 $\sin \alpha+\cos \alpha=\frac{1}{5}(0<\alpha<\pi)$, 求: \\
\begin{enumerate*}
    \item  $\sin \alpha-\cos \alpha$;
    \item  $\sin ^3 \alpha+\cos ^3 \alpha$;
    \item $\sin ^4 \alpha+\cos ^4 \alpha$.
\end{enumerate*}
\begin{jieda}
    \begin{enumerate}
        \item $(sin\alpha+cos\alpha)^2 = 1 + 2sin\alpha cos\alpha = \dfrac{1}{25}$，因此$sin\alpha cos\alpha = -\dfrac{12}{25}$ \\
        因此有$(sin\alpha - cos\alpha)^2 = 1 - 2sin\alpha cos\alpha = \dfrac{49}{25}$ \\
        故$sin\alpha - cos\alpha = \pm \dfrac{7}{5}$
        \\
        根据$\alpha \in (0, \pi)$，故$sin \alpha - cos\alpha > -1$，因此$sin\alpha - cos\alpha = \dfrac{7}{5}$ \\
        \ps 可以根据$sin\alpha + cos\alpha$在不同象限的值域直接判断出来$\alpha$是第二象限角.
        \item $sin^3 \alpha + cos^3 \alpha = (sin\alpha + cos\alpha)(sin^2 \alpha + cos^2 \alpha - sin\alpha cos\alpha) = (\dfrac{1}{5})(1 + \dfrac{12}{25}) = \dfrac{37}{125}$ \\
        \ps 也可以在第一问的结果上构建方程组$\left \{ \begin{array}{}
            sin\alpha + cos\alpha = \frac{1}{5} \\
            sin\alpha - cos\alpha = \frac{7}{5} 
        \end{array}\right.$解得$\left \{ \begin{array}{}
            sin\alpha = \frac{4}{5} \\
            cos\alpha = -\frac{3}{5} 
        \end{array}\right.$后带入计算（下同）
        \item $\sin ^4 \alpha+\cos ^4 \alpha = (\sin ^2 \alpha+\cos ^2 \alpha)^2 - 2(sin \alpha cos\alpha)^2 = 1 - 2 \times (-\dfrac{12}{25})^2 = \dfrac{337}{625}$
    \end{enumerate}
\end{jieda}
 
\end{enumerate}

C组：
\begin{enumerate}
    \item 求函数$f(x) = \sqrt{x^4 + 3x^2 -6x + 10} - \sqrt{x^4 - 3x^2 + 2x + 5}$的最大值.
    \begin{jieda}
        $f(x) = \sqrt{(x^2+1)^2 + (x-3)^2} - \sqrt{(x^2-2)^2 + (x+1)^2}$ \\
        其几何意义是点$P(x, x^2)$到$A(3, -1)$和$B(-1, 2)$的距离之差$|PA| - |PB|$.由三角形两边之和大于第三边可知$|PA| - |PB| \leq |AB|$，而等号可以取到.\\
        因此$f(x) \leq 5$，$x = \dfrac{\sqrt{89}-3}{8}$时取最大值.
    \end{jieda}
    \item 求函数$y = 2x + \sqrt{4x^2-8x+3}$的最小值.
    \begin{jieda}
        定义域先行，$4x^2-8x+3 \geq 0$，因此定义域为$(-\infty, \dfrac{1}{2}] \cup [\dfrac{3}{2}, +\infty)$ ，下面分析单调性
        \begin{enumerate}
            \item 当$x \geq \dfrac{3}{2}$时，容易判断函数单调递增
            \item 当$x \leq \dfrac{1}{2}$时，$f(x) = (2x-2) + \sqrt{(2x-2)^2-1}+2 = 2 + \dfrac{1}{(2x-2) - \sqrt{(2x-2)^2-1}}$，可以判断函数单调递减.
        \end{enumerate}
        因此最值应在端点处取得，$f(\dfrac{1}{2}) = 1, f(\dfrac{3}{2}) = 3$，因此函数最小值为$1$
    \end{jieda}
    \item 已知实数$a > b > 0$，求函数$f(x) = \dfrac{x}{\sqrt{a-x^2}-\sqrt{b-x^2}}$的最大值.
    \begin{jieda}
        $f(x) = \dfrac{x}{\sqrt{a-x^2}-\sqrt{b-x^2}} = \dfrac{x(\sqrt{a-x^2}+\sqrt{b-x^2})}{a-b} \\
        = \dfrac{1}{a-b}\sqrt{x^2(a+b)-2x^4+2x^2\sqrt{(a-x^2)(b-x^2)}} \\
        = \dfrac{1}{a-b}\sqrt{ab - [x^2 - \sqrt{(a-x^2)(b-x^2)}]^2} \\
        \leq \dfrac{\sqrt{ab}}{a-b}$
        当且仅当$x = \sqrt{\dfrac{ab}{a+b}}$时取等号.
    \end{jieda}
\end{enumerate}



\section{任意角的正弦、余弦、正切、余切（2）}


\begin{enumerate}
\item 已知$\tan\alpha=\dfrac{2ab}{a^2-b^2}$, 其中$a>b>0$, $\alpha\in(0,\frac{\pi}{2})$, 则$\sin\alpha$=\blkda{$\dfrac{2ab}{a^2+b^2}$}.
\begin{jieda}
$\sin\alpha=\dfrac{2ab}{\sqrt{(2ab)^2+(a^2-b^2)^2}}$. 
\end{jieda}


\item 已知$\dfrac{\cos\alpha}{1+\sin\alpha}=-\dfrac{1}{2}$, 则$\dfrac{\cos\alpha}{\sin\alpha-1}$=\blkda{2}. 
\begin{jieda}
结合$\cos^2\alpha+\sin^2\alpha=1$解得$(\cos\alpha,\sin\alpha)=(-\frac{4}{5},\frac{3}{5})$(另一组增根已舍去), 
代入所求即得. 
\end{jieda}



\item 若$\cot(\sin\theta)\cdot\tan(\cos\theta)>0$, 则$\theta$是第\blkda{I、III}象限角. 
\begin{jieda}
四个象限讨论一下即可. $\cos\theta,\sin\theta\in[-1,1]\subseteq(-\frac{\pi}{2},\frac{\pi}{2})$. 
\end{jieda}



\item 若$\theta$是锐角, $\sin\theta-\cos\theta=\dfrac{1}{2}$, 则$\sin^3\theta-\cos^3\theta$=\blkda{$\dfrac{11}{16}$}. 
\begin{jieda}
围绕$(\sin\theta\pm\cos\theta)^2=1\pm2\sin\theta\cos\theta$展开. 
\begin{align*}
\sin^3\theta-\cos^3\theta & =(\sin\theta-\cos\theta)(\sin^2\theta+\sin\theta\cos\theta+\cos^2\theta)\\
& = (\sin\theta-\cos\theta)\left(1+\dfrac{1-(\sin\theta-\cos\theta)^2}{2}\right)\\
& = \dfrac{1}{2}\cdot\dfrac{3-\frac{1}{4}}{2}\\
& = \dfrac{11}{16}.
\end{align*}
\end{jieda}



\item 设$\cot{x}=2$, 则$\dfrac{2\cos{x}-4\sin{x}}{5\cos{x}+3\sin{x}}$=\blkda{0}, 
$3\sin^2x-4\cos^2x$=\blkda{$-\dfrac{13}{5}$}. 
\begin{jieda}
$\dfrac{2\cos{x}-4\sin{x}}{5\cos{x}+3\sin{x}}=\dfrac{2\cot{x}-4}{5\cot{x}+3}=0$. \\
$3\sin^2x-4\cos^2x=\dfrac{3\sin^2x-4\cos^2x}{\cos^2x+\sin^2x}=\dfrac{3-4\cot^2x}{\cot^2x+1}=-\dfrac{13}{5}$. 
\end{jieda}

\item 用列举法写出集合$A=\left\{y\mid y=\dfrac{1}{\cos\alpha\sqrt{1+\tan^2\alpha}}
+\dfrac{2\tan\alpha}{\sqrt{\sec^2\alpha-1}}\right\}$=\blkda{$\{\pm1,\pm3\}$}.
\begin{jieda}
$y=\dfrac{\abs{\cos\alpha}}{\cos\alpha}+\dfrac{2\tan\alpha}{\abs{\tan\alpha}}$. 
\end{jieda}


\item 使函数$y=\sqrt{\sin{x}}+\sqrt{16-x^2}$有意义的$x$的取值范围是\blkda[3]{$[-4,-\pi]\cup[0,\pi]$}; 
使函数$y=\lg\sin(\cos{x})$有意义的$x$的取值范围是\blkda[5]{$\left(-\dfrac{\pi}{2}+2k\pi,\dfrac{\pi}{2}+2k\pi\right)(k\zz)$}. 
\begin{jieda}
$\sin{x}\geq 0\land 16-x^2\geq 0$, 故$x\in[-4,-\pi]\cup[0,\pi]$. \\
$\sin(\cos{x})>0\land\cos{x}\in[-1,1]\iff\cos{x}\in(0,1]\iff x\in\left(-\dfrac{\pi}{2}+2k\pi,\dfrac{\pi}{2}+2k\pi\right)(k\zz)$. 
\end{jieda}


\item 若$\tan\alpha=\sqrt3$, $\pi<\alpha<\dfrac{3\pi}{2}$, 那么$\cos\alpha-\sin\alpha$的值是\blkda{$\dfrac{-1+\sqrt3}{2}$}.
\begin{jieda}
由条件可取终边上一点$(-1,-\sqrt{3})$. 
\end{jieda}


\item 已知$\sin\alpha=m(\abs{m}<1)$, $\dfrac{\pi}{2}<\alpha<\dfrac{3\pi}{2}$, 那么$\tan\alpha=$\blkda{$-\dfrac{m}{\sqrt{1-m^2}}$}.
\begin{jieda}
因为$\dfrac{\pi}{2}<\alpha<\dfrac{3\pi}{2}$, 故终边上点的横坐标为负. \\
$m=\dfrac{m}{1}$, 故终边上取点$(-\sqrt{1-m^2},m)$, 于是立得. 
\end{jieda}


\item 若角$\alpha$的终边落在直线$x+y=0$上, 则$\dfrac{\sin\alpha}{\sqrt{1-\sin^2\alpha}}+\dfrac{\sqrt{1-\cos^2\alpha}}{\cos\alpha}$的值等于\blkda{0}.
\begin{jieda}
考虑一下符号基本上就可以把答案选出来了. 
\end{jieda}


\item 已知$\theta\in\left(0,\dfrac{\pi}{2}\right)$, 且$\log_{\tan\theta+\cot\theta}\sin\theta=-\dfrac{3}{4}$, 
则$\log_{\tan\theta}\cos\theta=$\blkda{$\dfrac{1}{2}$}.
\begin{jieda}
$(\tan\theta+\cot\theta)^{-\frac{3}{4}}=\sin\theta$, 
即$1=\sin^4\theta(\tan\theta+\cot\theta)^3=\sin^4\theta\cdot\left(\dfrac{\sin\theta}{\cos\theta}+\dfrac{\cos\theta}{\sin\theta}\right)^3$. \\
于是$1=\dfrac{\sin^4\theta}{\sin^3\theta\cos^3\theta}=\dfrac{\sin\theta}{\cos^3\theta}=\dfrac{\tan\theta}{\cos^2\theta}$. \\
故$\cos^2\theta=\tan\theta$, 所以$\log_{\tan\theta}\cos\theta=\dfrac{1}{2}$. 
\end{jieda}



\item 证明下列恒等式: 
\begin{enumerate}[(1)]
\item $\dfrac{1+2\sin{x}\cos{x}}{\cos^2x-\sin^2x}=\dfrac{1+\tan{x}}{1-\tan{x}}$; 
\item $\dfrac{1-\sin^6x-\cos^6x}{1-\sin^4x-\cos^4x}=\dfrac{3}{2}$;
\item $\dfrac{\tan\alpha\cdot\sin\alpha}{\tan\alpha-\sin\alpha}=\dfrac{\tan\alpha+\sin\alpha}{\tan\alpha\cdot\sin\alpha}$; 
\item $\dfrac{\cos\alpha}{1+\sin\alpha}-\dfrac{\sin\alpha}{1+\cos\alpha}=\dfrac{2(\cos\alpha-\sin\alpha)}{1+\sin\alpha+\cos\alpha}$. 
\end{enumerate}
\begin{jieda}
\begin{enumerate}[(1)]
\item $\LHS=\dfrac{(\sin{x}+\cos{x})^2}{(\cos{x}-\sin{x})(\cos{x}+\sin{x})}=\dfrac{\cos{x}+\sin{x}}{\cos{x}-\sin{x}}$. 

\item ``$\sin^2\theta+\cos^2\theta=1$''. 
\begin{align*}
\LHS & = \dfrac{1-(\sin^2x+\cos^2x)(\sin^4x-\sin^2x\cos^2x+\cos^4x)}{1-(\sin^2x+\cos^2x)^2+2\sin^2x\cos^2x}\\
& = \dfrac{1-(\sin^2x+\cos^2x)^2+3\sin^2x\cos^2x}{2\sin^2x\cos^2x}\\
& = \dfrac{3\sin^2x\cos^2x}{2\sin^2x\cos^2x}
\end{align*}

\item $\iff\tan^2\alpha\sin^2\alpha=\tan^2\alpha-\sin^2\alpha\iff\sin^2\alpha=\tan^2\alpha-\tan^2\alpha\sin^2\alpha$. 

\item 
\begin{align*}
\LHS & = \dfrac{1-sin\alpha}{cos\alpha} - \dfrac{1-cos\alpha}{sin\alpha} \\
& = \dfrac{sin\alpha - sin^2\alpha - (cos\alpha - cos^2\alpha)}{sin\alpha \cdot cos\alpha} \\
& = \dfrac{(cos\alpha - sin\alpha) \cdot (sin\alpha + cos\alpha - 1)}{sin\alpha \cdot cos\alpha} \\
& = \dfrac{cos\alpha - sin\alpha}{sin\alpha \cdot cos\alpha} \cdot \dfrac{(sin\alpha + cos\alpha)^2 - 1}{sin\alpha + cos\alpha + 1} \\
& = \dfrac{cos\alpha - sin\alpha}{sin\alpha \cdot cos\alpha} \cdot \dfrac{2sin\alpha cos\alpha}{sin\alpha + cos\alpha + 1} \\
& = \RHS
\end{align*}
\end{enumerate}
\end{jieda}


\item 已知$\cot\alpha=\dfrac{1}{3}$, 求$\sin\alpha$、$\cos\alpha$及$\tan\alpha$. 
\begin{jieda}
$\alpha$为第一象限角时, $\sin\alpha=\dfrac{3}{\sqrt{10}}$, $\cos\alpha=\dfrac{1}{\sqrt{10}}$, $\tan\alpha=3$; \\
$\alpha$为第三象限角时, $\sin\alpha=-\dfrac{3}{\sqrt{10}}$, $\cos\alpha=-\dfrac{1}{\sqrt{10}}$, $\tan\alpha=3$. 
\end{jieda}

\item 已知$\sin\alpha=\dfrac{m^2-1}{m^2+1}(m>0)$, 求$\cos\alpha$与$\tan\alpha$的值. 
\begin{jieda}
$\alpha$为第一或第四象限角时, $\cos\alpha=\dfrac{2m}{m^2+1}$, $\tan\alpha=\dfrac{m^2-1}{2m}$; \\
$\alpha$为第二或第三象限角时, $\cos\alpha=-\dfrac{2m}{m^2+1}$, $\tan\alpha=-\dfrac{m^2-1}{2m}$. 
\end{jieda}



\item 已知$\tan\alpha=\sqrt{2}$, 求下列各式的值:
\begin{enumerate}[(1)]
\item $\sin\alpha+2\cos\alpha$; 
\item $\dfrac{\cos\alpha-5\sin\alpha}{3\cos\alpha+\sin\alpha}$; 
\item $\dfrac{\sin^2\alpha-\sin\alpha\cos\alpha-3\cos^2\alpha}{5\sin\alpha\cos\alpha+\sin^2\alpha+1}$; 
\item $2\sin^2\alpha-\sin\alpha\cos\alpha+\cos^2\alpha$. 
\end{enumerate}
\begin{jieda}
齐次化. \\
\begin{enumerate*}[(1)]
\item $\pm\dfrac{\sqrt{6}+2\sqrt{3}}{3}$; 
\item $\dfrac{13-16\sqrt2}{7}$; 
\item $-\dfrac{1}{5}$; 
\item $\dfrac{5-\sqrt{2}}{3}$. 
\end{enumerate*}
\end{jieda}



\item
\begin{enumerate}
\item 已知$3\sin^2\alpha+2\sin^2\beta=5\sin\alpha$, 求$\sin^2\alpha+\sin^2\beta$的范围; 
\item 已知$6\sin3\alpha-\cos^22\beta=6$, 求角$\alpha$的值. 
\end{enumerate}
\begin{jieda}
\begin{enumerate}
\item $2\sin^2\beta=5\sin\alpha-3\sin^2\alpha\in[0,2]\land\sin\alpha\in[-1,1]$, 则$\sin\alpha\in[0,\frac{2}{3}]\cup\{1\}$. \\
故$E=\sin^2\alpha+\frac{5}{2}\sin\alpha-\frac{3}{2}\sin^2\alpha=-\frac{1}{2}\sin^2\alpha+\frac{5}{2}\sin\alpha\in[0,\frac{13}{9}]\cup\{2\}$. 

\item $6\sin3\alpha-6=\cos^22\beta$. \\
$\LHS\in[-12,0],\RHS\in[0,1]$, 故$6\sin3\alpha-6=\cos^22\beta=0$, \\
于是$3\alpha=2k\pi+\dfrac{\pi}{2}$, 即$\alpha=\dfrac{2k\pi}{3}+\dfrac{\pi}{6}(k\zz)$. 
\end{enumerate}
\end{jieda}

\end{enumerate}

C组：
\begin{enumerate}
    \item 设实数$x, y, z, w$满足$x+y+z+w = 1$，求$M = xw + 2yw + 3xy + 4xz + 5yz$的最大值.
    \begin{jieda}
        $w = 1 - (x+y+z)$，消去$w$有$M = x\cdot(1-x-y-z) + 2y \cdot (1-x-y-z) + 3xy  + 4xz + 5yz \\ = -x^2 + x - 2y^2 + 2y - 3z^2 + 3z = -(x-\dfrac{1}{2})^2 - 2(y - \dfrac{1}{2})^2 - 3(z - \dfrac{1}{2})^2 + \dfrac{3}{2}$ \\
        故当$x = y = z = \dfrac{1}{2}, w = -\dfrac{1}{2}, M$取得最大值$\dfrac{3}{2}$
    \end{jieda}
    \item 已知实数$m$满足：当关于$x$的实系数一元二次方程$ax^2+bx+c=0$有实数根时，$(a-b)^2+(b-c)^2+(c-a)^2 \geq ma^2$恒成立，求$m$的最大值.
    \begin{jieda}
        不等式变形为$\dfrac{(a-b)^2}{a^2}+\dfrac{(b-c)^2}{a^2}+\dfrac{(c-a)^2}{a^2} \geq m$，即$(1-\dfrac{b}{a})^2 + (\dfrac{b}{a} - \dfrac{c}{a})^2 + (\dfrac{c}{a} - 1)^2 \geq m$。设方程的两根为$x_1, x_2$，故根据韦达定理有$(1+x_1+x_2)^2 + (x_1+x_2+x_1x_2)^2 + (x_1x_2-1)^2 \geq m$。\\
        展开得$1+x_1^2+x_2^2+2x_1 + 2x_2+2x_1x_2 + x_1^2 + x_2^2 + x_1^2x_2^2 + 2x_1x_2 + 2x_1^2x_2 + 2x_1x_2^2 + x_1^2x_2^2 -2x_1x_2 + 1 \geq m$
        整理得$2(x_1^2x_2^2 + x_1^2x_2 + x_1x_2^2 + x_1^2 + x_2^2 + x_1x_2 + x_1 + x_2 + 1) \geq m$ \\
        观察式子可知$x_1, x_2$对称，共9项，故可尝试因式分解。得到$2(x_1^2 + x_1 + 1)(x_2^2 + x_2 + 1) \geq m$。可以看成以$x_1, x_2$为变量的两个二次函数相乘，该二次函数当自变量取$-\dfrac{1}{2}$时取最小值$\dfrac{3}{4}$，因此$m \leq 2 \times \dfrac{3}{4} \times \dfrac{3}{4} = \dfrac{9}{8}$
    \end{jieda}
    \item 设正数$x, y$满足$x^2+y^2+\dfrac{1}{x} + \dfrac{1}{y} = \dfrac{27}{4}$，求$P = \dfrac{15}{x} - \dfrac{3}{4y}$的最小值.
    \begin{jieda}
        $x^2+y^2+\dfrac{16}{x} + \dfrac{1}{4y} = \dfrac{27}{4} + P$ \\
        $\dfrac{27}{4} + P \geq \sqrt[3]{x^2 \cdot \dfrac{8}{x} \cdot \dfrac{8}{x}} + \sqrt[3]{y^2 \cdot \dfrac{1}{8y} \cdot \dfrac{1}{8y}} = 4 + \dfrac{1}{4} = \dfrac{25}{4}$
    \end{jieda}
    \item 记$F(x, y) = (x-y)^2+(\dfrac{x}{3} + \dfrac{3}{y})^2(y \neq 0)$，求$F(x, y)$的最小值.
    \begin{jieda}
        $F(x, y)$为点$A(x, \dfrac{x}{3})$到点$B(y, -\dfrac{3}{y})$的距离平方. \\
        由几何意义知应寻找双曲线上斜率为$\dfrac{1}{3}$的点，根据反比例函数的导数可得该点应为$(-3, 1)$或$(3, -1)$ \\
        根据两平行直线间的距离公式知$F(x, y)$的最小值为$(\dfrac{6}{\sqrt{10}})^2 = \dfrac{18}{5}$
    \end{jieda}
\end{enumerate}



\section{诱导公式（1）}
\zsd 奇变偶不变，符号看象限——偶

A组：
\begin{enumerate}
\item 若 $P(-3,4)$ 为角 $a$ 终边上一点, 则 $\cos (2 \pi-a)=$ \blkda{$-\dfrac{3}{5}$}.
\item 已知角 $\theta$ 的终边过点 $(1,-2)$, 则 $\cos (\pi+\theta)=$ \blkda{$-\dfrac{\sqrt{5}}{5}$}.
\item 若 $\cos \alpha=-\frac{3}{5}, a$ 是第二象限角, 则 $\sin (\alpha+\pi)=$ \blkda{$-\dfrac{4}{5}$}.
\item 已知 $1-\cos (\pi-\alpha)=2 \sin \alpha$, 则 $\tan \alpha=$\blkda{0或$\dfrac{3}{4}$}.
\item 下列等式中不恒成立的是\dotfill(\kh{D})
\xx{$\sin (2 \pi-\alpha)=-\sin \alpha$}
{$\cos (2 \pi-\alpha)=\cos \alpha$}
{$\tan (2 \pi-\alpha)=-\tan \alpha$}
{$\cot (2 \pi-\alpha)=\cot \alpha$}

\item 已知角 $\alpha$ 和 $\beta$ 的终边关于 $y$ 轴对称，下列等式恒成立的是\dotfill(\kh{A})
\xx{$\sin \alpha=\sin \beta$}
{$\cos \alpha=\cos \beta$}
{$\tan \alpha=\tan \beta$}
{$\cot \alpha=\cot \beta$}

\item 已知 $\sin \alpha=\frac{3}{5}$, 且 $\textbf{}$ 是第二象限角, 则 $\cos (\pi-\alpha)+\sin (\pi+\alpha)$ 的值等于\blkda{$\dfrac{1}{5}$}.

\item 函数 $y=\frac{|\sin \alpha|}{\sin \alpha}+\frac{2 \cos \alpha}{|\cos \alpha|}+\frac{|\tan \alpha|}{\tan \alpha}+\frac{2 \cot \alpha}{|\cot \alpha|}$ 的值域为\blkda{$\{-4, -2, 0, 6\}$}.


\end{enumerate}

B组：
\begin{enumerate}
\item 若$\cos(\pi-\alpha)=\dfrac{\sqrt5}{3}$, 且$\alpha\in\left(\dfrac{\pi}{2},\pi\right)$, 则$\sin(\pi+\alpha)$=\blkda{$-\dfrac{2}{3}$}. 

\item 已知$\sin(\alpha-\pi)=\dfrac{2}{3}$, 且$\alpha\in\left(-\dfrac{\pi}{2},0\right)$, 则$\tan\alpha$=\blkda{$-\dfrac{2}{\sqrt{5}}$}. 

\item 已知$\sin\left(\alpha-\dfrac{2\pi}{3}\right)=\dfrac{1}{4}$, 则$\sin\left(\alpha+\dfrac{\pi}{3}\right)$=\blkda{$-\dfrac{1}{4}$}. 


\item 已知$\cos\left(\dfrac{\pi}{6}+\theta\right)=\dfrac{\sqrt{3}}{3}$, 则$\cos\left(\dfrac{5\pi}{6}-\theta\right)$=\blkda{$\dfrac{-\sqrt{3}}{3}$}. 

\item 已知$\cos(11\pi-3)=p$, 用$p$表示$\tan(-3)$=\blkda{$\dfrac{\sqrt{1-p^2}}{p}$}.  


\item 若$\tan\alpha=-2$, 则$\sin(\alpha-\pi)\cdot\cos(\pi+\alpha)$=\blkda{$-\dfrac{2}{5}$}. 

\item $\sin^2(\alpha+\pi)-\cos(\alpha+\pi)\cos(-\alpha)+1$的值是\blkda{2}.

\item 在$\triangle{ABC}$中, 给出下列四个式子: 
\ding{172} $\sin(A+B)+\sin{C}$, 
\ding{173} $\cos(A+B)+\cos{C}$, 
\ding{174} $\sin(2A+2B)+\sin2C$, 
\ding{175} $\cos(2A+2B)+\cos2C$, 其中恒为常数的是\blkda{\ding{173}\ding{174}}.


\item 函数式$\sqrt{1+2\sin(\pi-2)\cos(\pi+2)}$化简的结果是\dotfill(\kh{A})
\xx{$\sin2-\cos2$}
{$\pm(\sin2-\cos2)$}
{$\cos2-\sin2$}
{以上结论都不对}

\item 当$n\zz$时, 在
\ding{172} $\sin(n\pi+\frac{\pi}{3})$; 
\ding{173} $\sin(2n\pi\pm\frac{\pi}{3})$; 
\ding{174} $\sin(n\pi+(-1)^n\frac{\pi}{3})$; 
\ding{175} $\cos(2n\pi\pm\frac{\pi}{6})$; 
\ding{176} $\cos(n\pi+(-1)^n\frac{\pi}{6})$中与$\sin\frac{\pi}{3}$相等的是\blkda{\ding{174}\ding{175}}.


\item 计算 $\cos 1^{\circ}+\cos 2^{\circ}+\cos 3^{\circ}+\cdots+\cos 177^{\circ}+\cos 178^{\circ}+\cos 179^{\circ}=$\blkda{0}.

\item 已知 $\cos 170^{\circ}=m$, 则 $\tan 10^{\circ}$ 的值为\blkda{$-\dfrac{\sqrt{1-m^2}}{m}$}.


\item 已知$\sin(-\pi+\alpha)+2\cos(3\pi-\alpha)=0$, 计算: 
\begin{enumerate}[(1)]
\item $\dfrac{2\sin\alpha-\cos\alpha}{\sin\alpha+3\cos\alpha}$; 
\item $\sin^2\alpha+\sin\alpha\cos\alpha-3\cos^2\alpha$. 
\end{enumerate}
\begin{jieda}
条件可化为$-\sin\alpha-2\cos\alpha=0$, 故$\tan\alpha=-2$. 
\begin{enumerate}[(1)]
\item $E=\dfrac{2\tan\alpha-1}{\tan\alpha+3}=-5$. 
\item $E=\dfrac{\sin^2\alpha+\sin\alpha\cos\alpha-3\cos^2\alpha}{\sin^2\alpha+\cos^2\alpha}=\dfrac{\tan^2\alpha+\tan\alpha-3}{\tan^2\alpha+1}=-\dfrac{1}{5}$
\end{enumerate}
\end{jieda}


\item 设$f(\theta)=\dfrac{2\cos^3\theta+\sin^2(2\pi-\theta)+\cos(-\theta)-3}{2+2\cos^2(\pi+\theta)+\cos(2\pi-\theta)}$, 求$f(\frac{\pi}{3})$的值. 
\begin{jieda}
$f(\theta)=\dfrac{2\cos^3\theta+1-\cos^2\theta+\cos\theta-3}{2\cos^2\theta+\cos\theta-2}=\cos\theta-1$. \\
故$f(\frac{\pi}{3})=-\dfrac{1}{2}$. 
\end{jieda}


\item 已知$\sin(3\pi-\alpha)=\sqrt2\sin(2\pi+\beta)$, $\sqrt3\cos(-\alpha)=-\sqrt2\cos(\pi+\beta)$, 
且$0<\alpha<\pi$, $0<\beta<\pi$, 求$\sin\alpha,\sin\beta$. 
\begin{jieda}
$\begin{cases}
\sin\alpha=\sqrt{2}\sin\beta\\
\sqrt{3}\cos\alpha=\sqrt{2}\cos\beta
\end{cases}$, 两式两边平方再相加可得$\sin^2\alpha+3\cos^2\alpha=2$, 故$\sin^2\alpha=\dfrac{1}{2}$. \\
又$\alpha\in(0,\pi)$, 故$\sin\alpha=\dfrac{\sqrt{2}}{2}$. 
代入第一式可得$\sin\beta=\dfrac{1}{2}$. 
% $\sin\alpha=\dfrac{\sqrt{2}}{2}$, $\sin\beta=\dfrac{1}{2}$. 
\end{jieda}

\item 已知$f(x)=\dfrac{\sin(n\pi-x)\cos(n\pi-x)}{\cos((n+1)\pi-x)}\cdot\tan(x-n\pi)\cot(n\pi-x)(n\zz)$, 求$f(\frac{7\pi}{6})$. 
\begin{jieda}
$n$为偶数时, $f(x)=-\sin{x}$; $n$为奇数时, $f(x)=\sin{x}$. \\
$f(\frac{7\pi}{6})=\begin{cases}
-\frac{1}{2}, & n\text{为奇数}\\
\frac{1}{2}, & n\text{为偶数}
\end{cases}$. 
\end{jieda}

\end{enumerate}
C组：
\begin{enumerate}
    \item 已知函数$y = |sinx + cosx + tanx + cotx + secx + cscx|$，求函数的最小值.
    \begin{jieda}
        $y = |sinx + cosx + \dfrac{1}{sinx \cdot cosx} + \dfrac{sinx + cosx}{sinx \cdot cosx}|$，\\
        令$t = sinx + cosx, t \in [-\sqrt{2}, -1) \cup (-1, 1) \cup (1, \sqrt{2}]$，故$y = |t + \dfrac{1}{\frac{t^2 - 1}{2}} + \dfrac{t}{\frac{t^2 - 1}{2}}| = |(t-1) + \dfrac{2}{t-1} + 1| \geq 2\sqrt{2} - 1$，当$t = 1-\sqrt{2}$时取等号.
    \end{jieda}
    \item 已知$x, y, z \in \mathbf{R^+}$，满足$x+y+z = xyz$，求函数$f(x,y,z) = x^2(yz-1)+y^2(zx-1)+z^2(xy-1)$的最小值.
    \begin{jieda}
        $f(x,y,z) = x^2(yz-1)+y^2(zx-1)+z^2(xy-1) = (x+y+z)(xyz)-(x^2+y^2+z^2) \\
        = (x+y+z)^2-(x^2+y^2+z^2) = 2(xy+yz+zx)$ \\
        由柯西不等式知$(xy+yz+zx)(z+x+y) \geq (3\sqrt{xyz})^2 = 9xyz = 9(x+y+z)$，故$2(xy+yz+zx) \geq 18$，当$x = y = z = \sqrt{3}$时取等.
    \end{jieda}
    
    \item 已知实数$x, y$满足：$17(x^2+y^2)-30xy-16 = 0$，求$\sqrt{16x^2 + 4y^2-16xy - 12x+6y + 9}$的最大值.
    \begin{jieda}
由 $17\left(x^2+y^2\right)-30 x y-16=0$ 可得, $(x+y)^2+16(x$ $-y)^2=16$,
而
$$
\begin{aligned}
&\sqrt{16 x^2+4 y^2-16 x y-12 x+6 y+9} \\
& =\sqrt{(4 x-2 y)^2-3(4 x-2 y)+9} \\
& =\sqrt{[3(x-y)+(x+y)]^2-3[3(x-y)+(x+y)]+9} .
\end{aligned}
$$


令 $x+y=4 \cos \theta, x-y=\sin \theta$,
则 $3(x-y)+(x+y)=3 \sin \theta+4 \cos \theta=5 \sin (\theta+\varphi)$,
其中, $\sin \varphi=\frac{4}{5}, \cos \varphi=\frac{3}{5}, 3(x-y)+(x+y) \in[-5,5]$,\\
故当 $3(x-y)+(x+y)=-5$ 时,
$\sqrt{16 x^2+4 y^2-16 x y-12 x+6 y+9}=\sqrt{(-5)^2-3(-5)+9} =7$

    \end{jieda}
    
\end{enumerate}

\section{诱导公式（2）}
\zsd 奇变偶不变，符号看象限——奇

A组：
\begin{enumerate}
\item 已知 $\sin \alpha=\frac{2 \sqrt{5}}{5}$, 则 $\cos \left(\alpha-\frac{\pi}{2}\right)=$\blkda{$\frac{2\sqrt{5}}{5}$}.

\item 若 $\sin a=\frac{1}{3}, \alpha\in\left(\frac{\pi}{2}, \pi\right)$, 则 $\sin \left(\alpha-\frac{3 \pi}{2}\right)$ 的值为\blkda{$-\frac{2\sqrt{2}}{3}$}.

\item 已知点 $A$ 的坐标为 $(3,4)$, 将 $O A$ 绕坐标原点 $O$ 顺时针旋转 $\frac{\pi}{2}$ 至 $O A^{\prime}$, 则点 $A^{\prime}$ 的坐标为\blkda{$(4,-3)$}.
\item 下列等式中恒成立的是\dotfill(\kh{A})
\xx{$\sin \left(\frac{3 \pi}{2}-\alpha\right)=-\cos a$}
{$\cos \left(\frac{3 \pi}{2}-\alpha\right)=\sin a$}
{$\tan \left(\frac{3 \pi}{2}-\alpha\right)=-\cot \alpha$}
{$\cot \left(\frac{3 \pi}{2}-\alpha\right)=\cot \alpha$}


\item 已知角 $\alpha$ 的终边经过点 $P(-3,-4)$, 则 $\cos \left(\frac{\pi}{2}+\alpha\right)$ 的值为\blkda{$\frac{4}{5}$}.


\end{enumerate}

B组：
\begin{enumerate}
\item 已知$\sin(\frac{\pi}{2}-\alpha)=\frac{3}{5}$, 则$\cos(\pi+\alpha)$的值为\blkda{$-\dfrac{3}{5}$}. 

\item 已知$\cos(\frac{\pi}{2}+\varphi)=\dfrac{\sqrt3}{2}$, 且$\abs{\varphi}<\dfrac{\pi}{2}$, 则$\tan\varphi$=\blkda{$-\sqrt{3}$}. 
\begin{jieda}
$-\sin\varphi=\dfrac{\sqrt{3}}{2}$, 故$\sin\varphi=-\dfrac{\sqrt{3}}{2}$. \\
又$\abs{\varphi}<\dfrac{\pi}{2}$, 故$-\dfrac{\pi}{2}<\varphi<0$, 于是$\tan\varphi=-\sqrt{3}$. 
\end{jieda}

\item 如果$\cos(\pi+\alpha)=-\dfrac{1}{3}$, 那么$\sin\left(\frac{3\pi}{2}-\alpha\right)$=\blkda{$-\dfrac{1}{3}$}. 
\begin{jieda}
$-\dfrac{1}{3}=\cos(\pi+\alpha)=-\cos\alpha$, 故$\cos\alpha=\dfrac{1}{3}$. \\
于是$E=-\cos\alpha=-\dfrac{1}{3}$. 
\end{jieda}

\item 已知函数$f(x)=\sqrt2\cos(x-\frac{\pi}{12})$, 若$\cos\theta=\frac{3}{5}$, $\theta\in(\frac{3\pi}{2},2\pi)$, 
则$f(\theta-\frac{5\pi}{12})$=\blkda{$-\frac{4\sqrt{2}}{5}$}. 
\begin{jieda}
$f(\theta-\frac{5\pi}{12})=\sqrt2\cos(\theta-\frac{5\pi}{12}-\frac{\pi}{12})=\sqrt{2}\sin\theta=-\frac{4\sqrt{2}}{5}$
\end{jieda}


\item 已知$\sin\beta=\dfrac{1}{3}$, $\sin(\alpha+\beta)=1$, 则$\sin(2\alpha+\beta)$=\blkda{$\dfrac{1}{3}$}. 
\begin{jieda}
$\sin(\alpha+\beta)=1\iff\alpha+\beta=2k\pi+\dfrac{\pi}{2}\iff\alpha=\dfrac{\pi}{2}-\beta+2k\pi(k\zz)$. \\
故$\sin(2\alpha+\beta)=\sin(\pi-2\beta+4k\pi+\beta)=\sin\beta=\dfrac{1}{3}$. 
\end{jieda}


\item 设$f(n)=\sin\left(\frac{n\pi}{4}+\alpha\right)$, 则$f(n)\cdot f(n+4)+f(n+2)\cdot f(n+6)$=\blkda{$-1$}. 
\begin{jieda}
令$\beta=\dfrac{n\pi}{4}+\alpha$, 故$E=\sin\beta\sin(\pi+\beta)+\sin(\frac{\pi}{2}+\beta)\sin(\pi+\frac{\pi}{2}+\beta)=-\sin^2\beta-\cos^2\beta=-1$. 
\end{jieda}

\item 已知$\sin(\frac{\pi}{2}+\theta)+3\cos(\pi-\theta)=\sin(-\theta)$, 则$\sin\theta\cos\theta+\cos^2\theta=$\dotfill(\kh{C})
\xx{$\dfrac{1}{5}$}
{$\dfrac{2}{5}$}
{$\dfrac{3}{5}$}
{$\dfrac{\sqrt5}{5}$}
\begin{jieda}
$\cos\theta-3\cos\theta=-\sin\theta$, 故$\tan\theta=2$. \\
$E=\dfrac{\sin\theta\cos\theta+\cos^2\theta}{\sin^2\theta+\cos^2\theta}=\dfrac{\tan\theta+1}{\tan^2\theta+1}=\dfrac{3}{5}$. 
\end{jieda}

\item 已知点$P(\sin(\pi+\theta),\sin(\frac{3\pi}{2}-\theta))$在第三象限, 则角$\theta$所在的象限是\dotfill(\kh{A})
\xx{第一象限}{第二象限}{第三象限}{第四象限}
\begin{jieda}
$\begin{cases}
\sin(\pi+\theta)<0\\
\sin(\frac{3\pi}{2}-\theta)<0
\end{cases}$, 
于是$\begin{cases}
\sin\theta>0\\
\cos\theta>0
\end{cases}$. 
\end{jieda}

%\dotfill($\dfrac{7\pi}{6}$)

\item 已知角$\alpha$的终边上一点的坐标为$(\sin\frac{4\pi}{3},\cos\frac{4\pi}{3})$, 则角$\alpha$的最小正值为\blkda{$\dfrac{7\pi}{6}$}.
\begin{jieda}
法一: 点$(-\frac{\sqrt{3}}{2},-\frac{1}{2})$, 故所求为$\dfrac{7\pi}{6}$. \\
法二: $f(\frac{\pi}{2}-\theta)=g(\theta)$. 
\end{jieda}

\item 已知角$\alpha$的顶点与原点$O$重合, 始边与$x$轴的正半轴重合, 将$\alpha$的终边按顺时针方向旋转$\frac{\pi}{2}$后, 
过点$P(\frac{3}{5},\frac{4}{5})$, 则$\cos\alpha$等于\dotfill(\kh{A})
\xx{$-\dfrac{4}{5}$}
{$\dfrac{4}{5}$}
{$-\dfrac{3}{5}$}
{$\dfrac{3}{5}$}
\begin{jieda}
$\beta=\alpha-\frac{\pi}{2}$, 则$\alpha=\beta+\frac{\pi}{2}$. 
故$\cos\alpha=\cos(\beta+\frac{\pi}{2})=-\sin\beta=-\frac{4}{5}$. 
\end{jieda}

\item 已知$f(x)=\dfrac{\cos(\frac{\pi}{2}+x)\cos(-x)\sin(\frac{3\pi}{2}-x)}{\sin(-\pi-x)\cos(2\pi-x)}$. 
\begin{enumerate}[(1)]
\item 化简$f(x)$; 
\item 若$x$是第三象限角, 且$\tan{x}=2$, 求$f(x)$的值. 
\end{enumerate}
\begin{jieda}
\begin{enumerate}[(1)]
\item $f(x)=\dfrac{-\sin{x}\cos{x}(-\cos{x})}{\sin{x}\cos{x}}=\cos{x}$. 
\item $f(x)=-\dfrac{1}{\sqrt{5}}$. 
\end{enumerate}
\end{jieda}

\item 已知$\sin\alpha$是方程$3x^2-10x-8=0$的根, 且$\alpha$为第三象限角, \\
求$\dfrac{\sin(\alpha+\frac{3\pi}{2})\sin(\frac{3\pi}{2}-\alpha)\tan^2(2\pi-\alpha)\tan(\pi-\alpha)}
{\cos(\frac{\pi}{2}-\alpha)\cos(\frac{\pi}{2}+\alpha)}$的值. 
\begin{jieda}
易得$\sin\alpha=-\dfrac{2}{3}$. 
故
\begin{align*}
E & = \dfrac{-\cos\alpha(-\cos\alpha)\tan^2\alpha(-\tan\alpha)}{\sin\alpha(-\sin\alpha)}\\
& = \tan\alpha\\
& = \dfrac{2}{\sqrt{5}}.
\end{align*}
\end{jieda}


\item 化简
\begin{enumerate}
    \item $\sin (21 \pi-\alpha)+\cos \left(\frac{9 \pi}{2}-\alpha\right)+\tan \left(\frac{9 \pi}{4}-\alpha\right)-\sin (-\alpha-19 \pi)-\cos \left(-\alpha-\frac{27}{2} \pi\right)-\tan \left(-\alpha-\frac{7 \pi}{4}\right)$.
    \item $\frac{\sin \left(31^{\circ}+\alpha\right)}{\tan \left(27^{\circ}+\alpha\right)} \cdot \frac{\tan \left(747^{\circ}+\alpha\right)}{\cos \left(36^{\circ}+\alpha\right)} \cdot \frac{\cos \left(1116^{\circ}+\alpha\right)}{\sin \left(751^{\circ}+\alpha\right)}$.
\end{enumerate}
\begin{jieda}
    \begin{enumerate}
\item $ y_s=\sin (\pi-\alpha)+\cos \left(\frac{\pi}{2}-\alpha\right)+\tan \left(\frac{\pi}{4}-\alpha\right)-\sin (\pi-\alpha)-\cos \left(\frac{\pi}{2}-\alpha\right)-\tan \left(\frac{\pi}{4}-\alpha\right)=0 $
\item $y_s=\frac{\sin \left(31^{\circ}+\alpha\right) \cdot \cos \left(27^{\circ}+\alpha\right)}{\sin \left(27^{\circ}+\alpha\right)} \cdot \frac{\sin \left(27^{\circ}+\alpha\right)}{\cos \left(36^{\circ}+\alpha\right) \cos \left(27^{\circ}+\alpha\right)} \cdot \frac{\cos \left(36^{\circ}+\alpha\right)}{\sin \left(31^{\circ}+\alpha\right)}=1$
\end{enumerate}
\end{jieda}



\item 已知 $\alpha$ 是第三象限角, $f(\alpha)=\frac{\sin \left(\alpha-\frac{\pi}{2}\right)+\cos \left(\frac{3 \pi}{2}+\alpha\right)}{\cos (-\alpha-2021 \pi)+\sin (2022 \pi-\alpha)}$.
\begin{enumerate}
    \item 若 $\alpha=-\frac{32 \pi}{3}$, 求 $f(\alpha)$ 的值;
    \item 若 $\tan \alpha=2$, 求 $f(\alpha)$ 的值.
\end{enumerate}
\begin{jieda}
$f_{(a)}=\frac{-\cos \alpha+\sin \alpha}{-\cos \alpha-\sin \alpha}$
    \begin{enumerate}
        \item $f(\alpha)=f\left(\frac{4}{3} \pi\right)=-2+\sqrt{3}$
        \item $f(\alpha)=\frac{-1+\tan \alpha}{-1-\tan \alpha}=-\frac{1}{3}$
    \end{enumerate}
\end{jieda}


\end{enumerate}

C组：
\begin{enumerate}
\item
\begin{enumerate}
\item 已知集合 $A=\left\{a \left\lvert\, a=k \pi+\frac{\pi}{6}\right., k \in \mathbb{Z}\right\}$.\\
\ding{172}是否存在 $B=[a, b]$, 使 $A \cap B=\left\{-\frac{11 \pi}{6},-\frac{5 \pi}{6}, \frac{\pi}{6}, \frac{7 \pi}{6}\right\}$ 成立? 如果存在, 求出 $a, b$的范围; 如果不存在, 说明理由.\\
\ding{173}是否存在 $B=[a, b]$, 使 $A \cap B$ 有且仅有 4 个元素? 如界存在, 求出 $b-a$ 的最大范围；如果不存在，说明遇由.
\item 所有能使 $\tan a=\tan 3$ 成立的 $a$ 组成集合 $A$, 请你写出一个集合 $B$, 使 $B \subseteq A$,且 $B$ 的元意有无限个。
\end{enumerate}
\begin{jieda}
    \begin{enumerate}
        \item \ding{172}  $a \in\left[-\frac{17}{6} \pi,-\frac{11}{6} \pi\right]$, $b \in\left[\frac{7}{6} \pi, \frac{13}{6} \pi\right)$ \\
        \ding{173} 存在满足条件的$B$, $b-a \in[3 \pi, 5 \pi)$
        \item $B=\{\alpha|\alpha=2k\pi+3,k\in\mathbb{Z}\}$
    \end{enumerate}
\end{jieda}

\item 当 $s$ 和 $t$ 取到所有实数时，求 $(s+7-|\cos t|)^2+(s-2|\sin t|)^2$ 的最小值.
\begin{jieda}
设点 $A(s+7, s)$, 点 $B(|\cos t|, 2|\sin t|)$, 则 $(s+7-|\cos t|)^2+$ $(s-2|\sin t|)^2=|A B|^2$.

注意到 $|\cos t| \geqslant 0,2|\sin t| \geqslant 0$, 不妨设 $t \in\left[0, \frac{\pi}{2}\right]$, 则点 $B$ 的坐标为 $(\cos t, 2 \sin t)$.
易知点 $A(s+7, s)$ 为直线 $x-y-7=0$ 上的动点, 点 $B(\cos t, 2 \sin t)$ 到直线 $x-y-7=0$的距离为
$$
d=\frac{|\cos t-2 \sin t-7|}{\sqrt{2}}=\frac{7+2 \sin t-\cos t}{\sqrt{2}} .
$$
又函数 $d=\frac{7+2 \sin t-\cos t}{\sqrt{2}}$ 在 $\left[0, \frac{\pi}{2}\right]$ 上单调递增, 故当 $t=0$ 时, $d$ 取得最小值 $\frac{6}{\sqrt{2}}$, 则 $|A B|^2$ 的最小值为 $\left(\frac{6}{\sqrt{2}}\right)^2=18$.
所以 $(s+7-|\cos t|)^2+(s-2|\sin t|)^2$ 的最小值为 18 .
\end{jieda}

\item 若 $a>0, b>0$, 求 $\min \left\{\max \left(a, b, \frac{1}{a^2}+\frac{1}{b^2}\right)\right\}$.
\begin{jieda}
设 $\left\{\max \left(a, b, \frac{1}{a^2}+\frac{1}{b^2}\right)\right\}=m$, 则 $a \leqslant m, b \leqslant m$, $\frac{1}{a^2}+\frac{1}{b^2} \leqslant m$, 因此 $m \geqslant \frac{2}{m^2}, m \geqslant \sqrt[3]{2}$.且当 $a=b=\sqrt[3]{2}$ 时, $m \geqslant \sqrt[3]{2}$ 所以 $m$ 的最小值为 $\sqrt[3]{2}$.所以 $\min \left\{\max \left(a, b, \frac{1}{a^2}+\frac{1}{b^2}\right)\right\}=\sqrt[3]{2}$.
\end{jieda}

\end{enumerate}


\section{求角}
\begin{enumerate}
\item 已知$\sin2x=-\dfrac{\sqrt2}{2}, x\in[-\pi,\pi]$, 则$x$=\blkda[3]{$-\dfrac{3\pi}{8},-\dfrac{\pi}{8},\dfrac{5\pi}{8},\dfrac{7\pi}{8}$}. 
\begin{jieda}
$2x\in[-2\pi,2\pi]$, 故$2x=-\dfrac{3\pi}{4},-\dfrac{\pi}{4},\dfrac{5\pi}{4},\dfrac{7\pi}{4}$, 
故$x=-\dfrac{3\pi}{8},-\dfrac{\pi}{8},\dfrac{5\pi}{8},\dfrac{7\pi}{8}$. 
\end{jieda}

\item 已知$\tan\left(2x+\dfrac{\pi}{6}\right)=-1$, 则$x$=\blkda[3]{$\dfrac{k\pi}{2}-\dfrac{5\pi}{24}(k\zz)$}. 
\begin{jieda}
$2x+\dfrac{\pi}{6}=k\pi-\dfrac{\pi}{4}$, 故$x=\dfrac{k\pi}{2}-\dfrac{5\pi}{24}(k\zz)$
\end{jieda}


\item 方程$2\sin^2x+3\cos{x}=0$, $x\in[-2\pi,2\pi]$的解集为\blkda[3]{$\left\{\pm\dfrac{4\pi}{3},\pm\dfrac{2\pi}{3}\right\}$} .
\begin{jieda}
$\iff 2-2\cos^2x+3\cos{x}=0$, 故$\cos{x}=-\dfrac{1}{2}$(2已舍去), 于是$x=2k\pi\pm\dfrac{2\pi}{3}(k\zz)$. \\
又$x\in[-2\pi,2\pi]$, 故$x=-\dfrac{4\pi}{3},-\dfrac{2\pi}{3},\dfrac{2\pi}{3},\dfrac{4\pi}{3}$. 
\end{jieda}

\item 已知$\sin\dfrac{\alpha}{2}=-\dfrac{\sqrt{3}}{2}$, 且$\alpha$是第二象限角, 则角$\alpha$=\blkda[3]{$4k\pi-\dfrac{4\pi}{3}(k\zz)$}. 
\begin{jieda}
$\dfrac{\alpha}{2}=2k\pi-\dfrac{\pi}{3}$或$2k\pi-\dfrac{2\pi}{3}$, 故$\alpha=4k\pi-\dfrac{2\pi}{3}$或$4k\pi-\dfrac{4\pi}{3}$. \\
又$\alpha$是第二象限角, 故$\alpha=4k\pi-\dfrac{4\pi}{3}(k\zz)$
\end{jieda}


\item 方程$\sin\left(x+\dfrac{\pi}{6}\right)-\cos{x}=0$的解集为\blkda[4]{$\left\{x\mid x=k\pi+\dfrac{\pi}{6},k\zz\right\}$}. 
\begin{jieda}
$\sin\alpha=\cos\beta\iff \sin\alpha=\sin\left(\dfrac{\pi}{2}-\beta\right)\iff
\alpha=\dfrac{\pi}{2}-\beta+2k\pi$或$\alpha=\dfrac{\pi}{2}+\beta+2k\pi$. 
\begin{itemize}
\item $x+\dfrac{\pi}{6}=\dfrac{\pi}{2}-x+2k\pi\iff x=k\pi+\dfrac{\pi}{6}(k\zz)$; 
\item $x+\dfrac{\pi}{6}=\dfrac{\pi}{2}+x+2k\pi\iff 2k\pi+\dfrac{\pi}{3}=0$, 无解. 
\end{itemize}
\end{jieda}


\item 已知集合$A=\left\{\alpha\mid 2\sin\alpha-1\geq 0\right\}$, $B=\left\{\alpha\mid\sqrt2\cos\alpha+1\geq 0\right\}$, \\
则$A\cap B$=\blkda[7]{$\left\{\alpha\mid 2k\pi+\dfrac{\pi}{6}\leq\alpha\leq 2k\pi+\dfrac{3\pi}{4},k\zz\right\}$}. 
\begin{jieda}
$\sin\alpha\geq\dfrac{1}{2}\iff 2k\pi+\dfrac{\pi}{6}\leq\alpha\leq 2k\pi+\dfrac{5\pi}{6}$, \\
$\cos\alpha\geq-\dfrac{1}{\sqrt{2}}\iff 2k\pi-\dfrac{3\pi}{4}\leq\alpha\leq 2k\pi+\dfrac{3\pi}{4}$. \\
于是$A\cap B=\left\{\alpha\mid 2k\pi+\dfrac{\pi}{6}\leq\alpha\leq 2k\pi+\dfrac{3\pi}{4},k\zz\right\}$. 
\end{jieda}


\item ``$\alpha=\dfrac{\pi}{6}+2k\pi(k\zz)$''是``$\cos2\alpha=\dfrac{1}{2}$''的\dotfill(\kh{A})
\xx{充分而不必要条件}{必要而不充分条件}{充分必要条件}{既不充分也不必要条件}


\item 若三角方程$\sin{x}=0$与$\sin2x=0$的解集分别为$E,F$, 则\dotfill(\kh{A})
\xx{$E\subsetneqq F$}{$E\supsetneqq F$}{$E=F$}{$E\cap F=\varnothing$}


\item 方程$\sqrt{1-\cos^2x}=\cos{x}$的解集是\dotfill(\kh{D})
\xx
{$\left\{x\mid x=k\pi\pm\dfrac{\pi}{4}, k\zz\right\}$}
{$\left\{x\mid x=k\pi+\dfrac{\pi}{4}, k\zz\right\}$} 
{$\left\{x\mid x=k\pi-\dfrac{\pi}{4}, k\zz\right\}$}
{$\left\{x\mid x=2k\pi\pm\dfrac{\pi}{4}, k\zz\right\}$}


\item 已知$\theta\in[0,2\pi]$, $\sin\theta$和$\cos\theta$是方程$x^2-kx+k+1=0$的两个根, 求$\theta$. 
\begin{jieda}
$\Delta\geq 0$, 则$k\in(-\infty,2-2\sqrt{2}]\cup[2+2\sqrt{2},+\infty)$. \\
于是$\begin{cases}
\sin\theta+\cos\theta=k\\
\sin\theta\cos\theta=k+1
\end{cases}$, 
而$(\sin\theta+\cos\theta)^2=1+2\sin\theta\cos\theta$, \\
故$k^2=1+2(k+1)$, 解得$k=-1$或3. \\
当$k=3$时, $\sin\theta\cos\theta=4>1$, 不符合条件, 故舍去; \\
当$k=-1$时, 解得$(\cos\theta,\sin\theta)=(0,-1)$或$(-1,0)$. 又$\theta\in[0,2\pi]$, 故$\theta=\dfrac{3\pi}{2},\pi$. 
\end{jieda}


\item 已知集合$A=\left\{-\dfrac{1}{2},3,\cos{x}\right\}$, $B=\left\{\dfrac{\sqrt{3}}{2},\sin{x}\right\}$. 
\begin{enumerate}[(1)]
\item 当$B\subsetneqq A$时, 求$x$的值;
\item 当$A\cap B=\{y\}$时, 求$x$和$y$的值. 
\end{enumerate}
\begin{jieda}
\begin{enumerate}[(1)]
\item 由题意可得$\begin{cases}
-\frac{1}{2}=\sin{x}\\
\cos{x}=\frac{\sqrt{3}}{2}
\end{cases}$, 故$x=2k\pi+\dfrac{11\pi}{6}(k\zz)$. 

\item 若$y=\dfrac{\sqrt{3}}{2}$, 则$\cos{x}=\dfrac{\sqrt{3}}{2}\land\sin{x}\neq-\dfrac{1}{2}$, 故$x=2k\pi+\dfrac{\pi}{6}(k\zz)$. \\
若$y=\sin{x}$, 则$y=\cos{x}$或$y=-\dfrac{1}{2}$, 即 (i)$x=k\pi+\dfrac{\pi}{4}$, 或(ii)$x=2k\pi+\dfrac{7\pi}{6}$或$x=2k\pi+\dfrac{11\pi}{6}$. \\
检验可得$x=2k\pi+\dfrac{11\pi}{6}$不符合条件. \\
综上, $(x,y)=(2k\pi+\frac{\pi}{6},\frac{\sqrt{3}}{2})$, $(2k\pi+\frac{\pi}{4},\frac{\sqrt{2}}{2})$, 
$(2k\pi+\frac{5\pi}{4},-\frac{\sqrt{2}}{2})$, $(2k\pi+\frac{7\pi}{6},-\frac{1}{2})$, 其中$k\zz$. 
\end{enumerate}
\end{jieda}

\item 已知关于$x$的方程$\sin^2x+\cos{x}+m=0$. 
\begin{enumerate}
\item 当$m=1$时, 求方程的解; 
\item 要使此方程有解, 试确定$m$的取值范围. 
\end{enumerate}
\begin{jieda}
\begin{enumerate}[(1)]
\item $1-\cos^2x+\cos{x}+1=0$, 解得$\cos{x}=-1$(2已舍去), 故$x=2k\pi+\pi(k\zz)$. 

\item $m=-\sin^2x-\cos{x}=\cos^2x-\cos{x}-1\in[-\frac{5}{4},1]$. 
\end{enumerate}
\end{jieda}


\item 已知角 $a$ 是第三象限角, $\tan a=\frac{1}{2}$.
\begin{enumerate}
    \item 求 $\sin a, \cos a$ 的值;
    \item 求 $\frac{1+2 \sin (\pi-a) \cos (-2 \pi-a)}{\sin ^2(-a)-\sin ^2\left(\frac{5 \pi}{2}-a\right)}$ 的值.
\end{enumerate}
\begin{jieda}
\begin{enumerate}
    \item $\tan \alpha=\frac{\sin \alpha}{\cos \alpha}=\frac{1}{2}, \sin ^2 \alpha+\cos ^2 \alpha=1$,

所以 $\left\{\begin{array}{l}\sin \alpha=\frac{\sqrt{5}}{5} \\ \cos \alpha=\frac{2 \sqrt{5}}{5}\end{array}\right.$ 或 $\left\{\begin{array}{l}\sin \alpha=-\frac{\sqrt{5}}{5} \\ \cos \alpha=-\frac{2 \sqrt{5}}{5}\end{array}\right.$ ，
而角 $\alpha$ 是第三象限角，则 $\sin \alpha<0, \cos \alpha<0$ ，故 $\left\{\begin{array}{l}\sin \alpha=-\frac{\sqrt{5}}{5} \\ \cos \alpha=-\frac{2 \sqrt{5}}{5}\end{array}\right.$

\item $\frac{1+2 \sin (\pi-\alpha) \cos (-2 \pi-\alpha)}{\sin ^2(-\alpha)-\sin ^2\left(\frac{5 \pi}{2}-\alpha\right)}=\frac{1+2 \sin \alpha \cos (-\alpha)}{(-\sin \alpha)^2-\sin ^2\left(\frac{\pi}{2}-\alpha\right)}$

$$
\begin{aligned}
& =\frac{1+2 \sin \alpha \cos \alpha}{\sin ^2 \alpha-\cos ^2 \alpha}=\frac{\sin ^2 \alpha+\cos ^2 \alpha+2 \sin \alpha \cos \alpha}{\sin ^2 \alpha-\cos ^2 \alpha} \\
& =\frac{(\sin \alpha+\cos \alpha)^2}{(\sin \alpha+\cos \alpha)(\sin \alpha-\cos \alpha)}=\frac{\sin \alpha+\cos \alpha}{\sin \alpha-\cos \alpha}=\frac{\tan \alpha+1}{\tan \alpha-1} .
\end{aligned}
$$
因为 $\tan \alpha=\frac{1}{2}$ ，所以原式 $=\frac{\frac{1}{2}+1}{\frac{1}{2}-1}=-3$.
\end{enumerate}
\end{jieda}


\end{enumerate}

C组：
\begin{enumerate}
    \item 已知函数 $f(x)$ 对于任意实数 $x, y$ 都有 $f(x+y)=f(x)+f(y)-3$, 当 $x>0$时 $f(x)<3$
    \begin{enumerate}
        \item $f(x)$ 在实数集 $R$ 上是否为单调函数? 并说明理由.
        \item 若 $f(6)=-9$, 求 $f\left(\left(\frac{1}{2}\right)^{2010}\right)$.
    \end{enumerate}
    \begin{jieda}
    \begin{enumerate}
        \item 对任意的 $x_1 、 x_2 \in R$ ，且 $x_1<x_2$ ，则 $f\left(x_2\right)=f\left(x_1+\left(x_2-x_1\right)\right)=f\left(x_1\right)+f\left(x_2-x_1\right)-3$.
因 $x_2-x_1>0$ ，所以， $f\left(x_2-x_1\right)<3$ ，即 $f\left(x_2\right)<f\left(x_1\right)$.
从而， $f(x)$ 在 $R$ 上为单调减函数.
\item 由 $f(6)=f(2)+f(4)-3=f(2)+(f(2)+f(2)-3)-3=3 f(2)-6=-9$得 $f(2)=-1$. 又由 $f(2)=f(1)+f(1)-3$, 得 $f(1)=1$.
对任意的 $a \in R$, 显然有 $f(a)=2 f\left(\frac{a}{2}\right)-3 \Rightarrow f\left(\frac{a}{2}\right)=\frac{f(a)+3}{2}$.
令 $a_1=f(1)=1, a_{n+1}=f\left(\left(\frac{1}{2}\right)^n\right)=\frac{f\left(\left(\frac{1}{2}\right)^{n-1}\right)+3}{2}=\frac{a_n+3}{2}$.
令 $b_n=a_{n+1}-a_n$. 则 $b_n=\frac{1}{2} b_{n-1}, b_1=a_2-a_1=1$.
从而， $b_n=\left(\frac{1}{2}\right)^{n-1}$ ，即 $a_{n+1}-a_n=\left(\frac{1}{2}\right)^{n-1}$.
故 $a_{n+1}=a_1+1+\frac{1}{2}+\cdots+\left(\frac{1}{2}\right)^{n-1}=3-\left(\frac{1}{2}\right)^{n-1}$.
所以， $f\left(\left(\frac{1}{2}\right)^{2010}\right)=a_{2011}=3-\left(\frac{1}{2}\right)^{2009}$.
    \end{enumerate}
    \end{jieda}

\item 设 $a, b$ 为正实数, $\frac{1}{a}+\frac{1}{b} \leq 2 \sqrt{2},(a-b)^2=4 a^3 b^3$, 求 $\log _{\mathrm{a}} b$.
\begin{jieda}
由 $(a-b)^2=4(a b)^3$ 得$(a+b)^2=4 a b+4(a b)^3 \geqslant 8(a b)^2,$
故 $a+b \geqslant 2 \sqrt{2} a b$,
又因为 $\frac{1}{a}+\frac{1}{b} \leqslant 2 \sqrt{2}$, 则 $a+b \leqslant 2 \sqrt{2} a b$,
所以 $a+b=2 \sqrt{2} a b$,
又 $(a+b)^2=4 a b+4(a b)^3$,

令 $a b=t$, 则 $8 t^2=4 t+4 t^3$, 得 $t-1=0, a b=1, \log _a b=-1$.

\end{jieda}
\end{enumerate}